\chapter*{概率}
\subsection*{随机游走}
\begin{enumerate}
    \item 棱长为1的正四面体$A-BCD$,有一小虫从顶点A处开始按以下规则爬行: 在每一顶点处以同样的概率选择通过这个顶点的3条棱之一,并一直爬到这条棱的尽头,记小虫爬了$n$米后重新回到点A的概率为$P_n$.
    \begin{enumerate}
        \item 求$P_1$和$P_2$的值;
        \item 探寻$P_n$与$P_{n-1}$的关系;
        \item 求$P_n$的表达式.
    \end{enumerate}

    \begin{jieda}
        \begin{enumerate}
            \item $P_1=0$,$P_2=\dfrac{1}{3}$
            \item $P_n=\dfrac{1}{3}(1-P_{n-1})$
            \item $P_{n}=\dfrac{1}{4}+\dfrac{3}{4}(-\dfrac{1}{3})^n$
        \end{enumerate}
    \end{jieda}

    \item “布朗运动”是指微小颗粒永不停息的无规则随机运动,在如图所示的试验容器中,容器由三个仓组成,某粒子作布朗运动时每次会从所在仓的通道口中随机选择一个到达相邻仓或者容器外,一旦粒子到达容器外就会被外部捕获装置所捕获,此时试验结束. 已知该粒子初始位置在 1 号仓,则试验结束时该粒子是从 1 号仓到达容器外的概率为\blkda{$\dfrac{10}{13}$}
    \begin{center}
    \includegraphics[width=0.2\textwidth]{bodyxb/src/bulangyundong.png}
    \end{center}
    \begin{jieda}
        \begin{dingautolist}{172}
            \item 方法1
            
            不难发现，若粒子在仓内，则奇数时刻粒子仅可能在奇数仓，偶数时刻粒子仅可能在偶数仓，故将每两次移动并为一组考虑，设$P_n^1$为$n$组后（$n \in \mathbf{N}$）粒子在$1$号仓的概率，$P_n^3$为$n$组后（$n \in \mathbf{N}$）粒子在$3$号仓的概率.

            则有$\left\{\begin{array}{l}
                P_n^1 = \dfrac{1}{9}P_{n-1}^1 + \dfrac{1}{6}P_{n-1}^3 \\
                P_n^3 = \dfrac{1}{9}P_{n-1}^1 + \dfrac{1}{6}P_{n-1}^3 
            \end{array}\right.$，$\left\{\begin{array}{l}
                P_0^1 = 1 \\
                P_0^3 = 0
            \end{array}\right.$，故$n \geq 1$时$P_n^1 = P_n^3 = \dfrac{1}{9} \times \left(\dfrac{5}{18}\right)^{n-1}$

            设$P_n$为$n$组后从1 号仓到达容器外的概率，则$P_n = \dfrac{2}{3}P_n^1$，故$\mathop{\sum }\limits_{{n = 0}}^{{+\infty }}P_n = \dfrac{2}{3} \times \mathop{\sum }\limits_{{n = 0}}^{{+\infty }}P_n^1 = \dfrac{2}{3} \times \left(1+\dfrac{2}{13}\right) = \dfrac{10}{13}$.

            \item 方法2
            
            设从 $i$ 出发最终从 1 号口出的概率为 ${P}_{i}$ ,所以 $\left\{  \begin{array}{l} {P}_{1} = \dfrac{2}{3} + \dfrac{1}{3}{P}_{2} \\  {P}_{2} = \dfrac{1}{3}{P}_{1} + 0 + \dfrac{1}{3}{P}_{3} = \dfrac{1}{3}{P}_{1} + \dfrac{1}{6}{P}_{2} \\  {P}_{3} = \dfrac{1}{2}{P}_{2} \end{array}\right.$，解得${P}_{1} = \dfrac{10}{13}$.
        \end{dingautolist}
    \end{jieda}
\end{enumerate}